% =====================================================================
%  phuluc_mcp.tex  —  PHỤ LỤC B: PROTOTYPE ĐIỀU PHỐI BẰNG MCP
% =====================================================================

\chapter*{PHỤ LỤC B. PROTOTYPE ĐIỀU PHỐI BẰNG MCP}
\addcontentsline{toc}{chapter}{PHỤ LỤC B. PROTOTYPE ĐIỀU PHỐI BẰNG MCP}
\label{app:mcp}

Model Context Protocol (MCP) là giao thức chuẩn hóa cho phép mô hình ngôn ngữ
lớn (LLM) gọi công cụ bên ngoài theo mô hình máy chủ--máy khách. Phụ lục này mô
tả một prototype dùng MCP làm lớp giao tiếp phía trên engine IPPO của khóa luận.
Prototype không tham gia trả lời các câu hỏi nghiên cứu.

\section*{B.1. So sánh hai kiến trúc điều phối}

\begin{table}[htbp]
    \centering
    \caption{So sánh kiến trúc IPPO và MCP theo các chiều điều phối}
    \label{tab:soSanhMCP}
    \small
    \begin{tabular}{lp{5.0cm}p{5.0cm}}
        \toprule
        \textbf{Khía cạnh} & \textbf{IPPO (đề tài này)} & \textbf{Kiến trúc MCP} \\
        \midrule
        Phân công nhiệm vụ & Không có: mỗi tác tử có nhiệm vụ cố định (một cặp kho--mặt hàng) & Host/LLM phân công tường minh cho từng máy chủ \\
        Chọn hành động & Mạng chính sách $\pi_\theta(a \mid o)$ ánh xạ quan sát sang hành động & LLM suy luận và chọn công cụ phù hợp \\
        Lấy ngữ cảnh & Vector quan sát $43$ chiều từ môi trường & Máy chủ cung cấp ngữ cảnh qua giao thức chuẩn \\
        Kết quả trả về & Phần thưởng và trạng thái kế tiếp từ môi trường & Kết quả công cụ trả về LLM để tiếp tục suy luận \\
        \bottomrule
    \end{tabular}
\end{table}

Hai kiến trúc phục vụ hai lớp bài toán khác nhau. IPPO phù hợp khi hành động là
lượng đặt hàng và tín hiệu phản hồi là chi phí. MCP phù hợp khi hệ thống cần
giao tiếp ngôn ngữ với người dùng và gọi nhiều công cụ đa dạng.

\section*{B.2. Luồng điều phối và công cụ}

Theo luồng ngữ cảnh $\rightarrow$ nhiệm vụ $\rightarrow$ chọn công cụ
$\rightarrow$ thực thi $\rightarrow$ trả kết quả: người quản lý kho đặt câu hỏi
(ví dụ ``tuần tới kho CA\_3 có mặt hàng nào sắp thiếu không?''); LLM nhận nhiệm
vụ và chọn công cụ; máy chủ MCP lấy ngữ cảnh từ môi trường mô phỏng (tồn kho,
hàng đang về, lịch) và gọi chính sách IPPO hoặc baseline; kết quả trả về LLM để
tổng hợp câu trả lời. IPPO vẫn là nơi ra quyết định định lượng.

Prototype (\texttt{app/mcp\_server.py}, Python MCP SDK) cung cấp năm công cụ và
một tài nguyên (Bảng~\ref{tab:congCuMCP}), tái sử dụng phần lõi mô phỏng của ứng
dụng Streamlit. Prototype đã được kiểm chứng bằng một client MCP chuẩn qua giao
thức stdio: client liệt kê được đủ công cụ và gọi thành công. Prototype chưa
được đánh giá với người dùng thật và chưa xử lý xác thực.

\begin{table}[htbp]
    \centering
    \caption{Công cụ của prototype MCP}
    \label{tab:congCuMCP}
    \small
    \begin{tabular}{lp{9.2cm}}
        \toprule
        \textbf{Công cụ / tài nguyên} & \textbf{Chức năng} \\
        \midrule
        \texttt{danh\_sach\_kho\_va\_mat\_hang} & Liệt kê $10$ kho, $30$ mặt hàng và chỉ số tương ứng \\
        \texttt{canh\_bao\_thieu\_hang} & Mô phỏng tới một ngày, trả về các cặp có tồn kho cộng hàng đang về dưới ngưỡng số ngày cầu \\
        \texttt{de\_xuat\_dat\_hang} & Tình trạng một cặp và lượng đặt đề xuất của IPPO cùng ba baseline \\
        \texttt{mo\_phong} & Chạy một chính sách trên toàn mạng, trả về chi phí, fill rate, số ngày--cặp thiếu hàng \\
        \texttt{what\_if} & So sánh các chính sách khi thay đổi nhu cầu, sức chứa hoặc thời gian giao \\
        \texttt{ket-qua://tong-hop} & Tài nguyên chứa kết quả đánh giá chính thức để LLM trích dẫn \\
        \bottomrule
    \end{tabular}
\end{table}
